% Part: incompleteness
% Chapter: incompleteness-provability
% Section: godels-paper

\documentclass[../../../include/open-logic-section]{subfiles}

\begin{document}

\olfileid[hi]{inc}{inp}{gop}

\olsection{ग्योडेल (G\"odel) के मूल शोधपत्र से तुलना}

ग्योडेल (G\"odel) के 1931 के शोधपत्र पर कुछ समय लगाना सार्थक है। उसका
परिचय उन विचारों की रूपरेखा देता है जिनकी हमने अभी चर्चा की। शोधपत्र को
केवल सरसरी तौर पर पढ़ने पर भी प्रत्येक चरण में चल रही बात समझना आसान है:
पहले ग्योडेल (G\"odel) औपचारिक तंत्र $P$ (वाक्यविन्यास, अभिगृहीत, प्रमाण-नियम) का
वर्णन करते हैं; फिर आदिम पुनरावर्ती फलन और संबंध परिभाषित करते हैं; फिर
दिखाते हैं कि $x B y$ आदिम पुनरावर्ती है और तर्क देते हैं कि आदिम
पुनरावर्ती फलनों तथा संबंधों का $\Th{P}$ में निरूपण होता है। इसके बाद वे
ऊपर की भाँति अपूर्णता-प्रमेय सिद्ध करते हैं। अनुभाग 3 में वे दिखाते हैं कि
असिद्ध कथन को अंकगणित की भाषा का वाक्य चुना जा सकता है। यही $\beta$-उपपत्ति
का उद्गम है; संगणनीय फलनों की~$\Th{Q}$ में निरूपणीयता दिखाते समय अनुक्रमों
को सँभालने के लिए हमने भी उसी का उपयोग किया। ग्योडेल (G\"odel) पर्याप्त अभिगृहीतों
का कोई न्यूनतम समुच्चय अलग करने तक नहीं जाते, किंतु अब हम जानते हैं कि
$\Th{Q}$ काम कर देता है। अंततः अनुभाग 4 में वे द्वितीय अपूर्णता-प्रमेय के
प्रमाण की रूपरेखा देते हैं।

\end{document}
